Emergency department overcrowding is worsened by low-acuity patients, and combining self-service kiosks with telemedicine may offer these patients an alternative channel inside the health unit itself. This work proposes and analytically assesses a software architecture for kiosk-based telemedicine, aimed at post-triage care of patients classified as green or blue by the Manchester Triage System in Brazilian public health units. The research is applied and follows the Design Science Research paradigm, combining a protocol-based narrative review, traceable requirements derivation, layered modelling with the C4 model, and ex ante analytical evaluation through quality scenarios, regulatory conformance checking and requirement-component coverage analysis. The result is a five-layer architecture: kiosk, teleconsultation platform, clinical backend, HL7 FHIR R4 gateway to the National Health Data Network and the primary care record system, and cross-cutting security. The proposal includes a traceability matrix with 19 requirements, clinical criteria for eligibility, reassessment and return to in-person care, a mapping of the sections of the national Clinical Encounter Record, a specification of authentication, validation and error handling based on the network's official documentation, a distinction between the legal basis for data processing and clinical consent, and privacy, contingency and accessibility controls. In the evaluation, all requirements were linked to components, 25 of 39 conformance items were met at design level, 9 were partially met and 5 could not be verified without a prototype or certification testing, and seven of nine scenarios were met at design level. This is a preliminary proposal, neither prototyped nor empirically evaluated. Its compatibility is design-level and depends on expert validation, workflow simulation, certification testing and cost and performance analysis.
